%% toc: P15 159 159.8
\chapter{Conditional Independence of Sigma-Algebras and Its Inference Rules}
\label{sec:p15-w02c-p15-009-sigma-ci}

This chapter states conditional independence at the level of
sub-sigma-algebras of an arbitrary probability space and proves the rules
by which conditional independence statements are inferred from one another.
No standard Borel hypothesis, density or positivity assumption is used.
The random-element form for standard Borel spaces, with its kernel
characterization, is \cref{def:p16-measurable-ci} and
\cref{thm:p16-measurable-ci-calculus}; item~7 of
\cref{thm:p15-w02c-p15-009-sigma-ci-criteria} shows that the two
definitions agree where both apply. The presentation follows Lauritzen's
chapter on conditional independence \cite[ch.~2, pp.~39--52]{Lauritzen2019GraphicalLectures2020}.

\emph{Inputs.} Conditional expectation and conditional probability given a
sub-sigma-algebra (\cref{def:p15-general-conditional-expectation}) with
existence and almost-sure uniqueness (\cref{thm:p15-conditional-existence}),
linearity and the known-information rule (\cref{thm:p15-conditional-basic}),
conditional monotone convergence (\cref{thm:p15-conditional-limits}), the
tower rule (\cref{thm:p15-conditional-tower}), taking out a known factor
(\cref{thm:p15-conditional-pullout}), finite partitions
(\cref{thm:p15-conditional-partitions}), and the pi-lambda theorem
(\cref{thm:p15-measure-uniqueness}).
\emph{Conclusions.} Criteria for conditional independence (generators,
products, the one-sided prediction form), the inference rules symmetry,
reduction, decomposition, weak union and contraction, and two exact
limitations: conditioning can create or destroy independence, and the
intersection rule fails without an extra hypothesis.

\begin{hbdefinition}{def:p15-w02c-p15-009-sigma-ci}{Conditional independence given information}
Fix a probability space \((\Omega,\mathcal F_0,P)\) and a
sub-sigma-algebra \(\mathcal H\subseteq\mathcal F_0\). Events
\(A,B\in\mathcal F_0\) are \emph{conditionally independent given
\(\mathcal H\)}, written \(A\perp B\mid\mathcal H\), when
\[
 P(A\cap B\mid\mathcal H)=P(A\mid\mathcal H)\,P(B\mid\mathcal H)
 \qquad\text{almost surely.}
\]
Classes of events \(\mathcal A,\mathcal B\subseteq\mathcal F_0\) are
conditionally independent given \(\mathcal H\), written
\(\mathcal A\perp\mathcal B\mid\mathcal H\), when \(A\perp B\mid\mathcal H\)
for every \(A\in\mathcal A\) and \(B\in\mathcal B\). For sub-sigma-algebras
\(\mathcal B,\mathcal H\), the sigma-algebra \(\mathcal B\vee\mathcal H\) is
the smallest one containing both. For random elements \(X,Y,Z\) with values
in measurable spaces, \(X\perp Y\mid Z\) means
\(\sigma(X)\perp\sigma(Y)\mid\sigma(Z)\), and a tuple of random elements
generates the sigma-algebra of the tuple.
\index{conditional independence!given a sigma-algebra}
\end{hbdefinition}

Lauritzen gives these definitions for events and for classes of events
\cite[Definitions~2.1 and~2.5, pp.~40 and~42]{Lauritzen2019GraphicalLectures2020}
and identifies a random variable with the sigma-algebra it generates in
such statements (p.~46). He writes the relation with a doubled
perpendicular sign; this handbook writes \(\perp\), as in
\cref{def:p16-measurable-ci}. Lauritzen stresses that the conditioning
sigma-algebra is part of the relation: refining or coarsening it can create
or destroy conditional independence (p.~40).
\Cref{thm:p15-w02c-p15-009-ci-tables} gives exact instances of both.

\begin{hbtheorem}{thm:p15-w02c-p15-009-sigma-ci-criteria}{proved=annex:p15-w02c-p15-009-sigma-ci-criteria}{Criteria for conditional independence}
Let \(\mathcal A,\mathcal B,\mathcal H\) be sub-sigma-algebras of
\(\mathcal F_0\).
\begin{enumerate}
\item If \(\mathcal C\) and \(\mathcal D\) are pi-systems with
      \(\mathcal C\perp\mathcal D\mid\mathcal H\), then
      \(\sigma(\mathcal C)\perp\sigma(\mathcal D)\mid\mathcal H\).
\item If \(\mathcal A\perp\mathcal B\mid\mathcal H\), \(X\) is an
      \(\mathcal A\)-measurable and \(Y\) a \(\mathcal B\)-measurable real
      random variable, and \(E|X|\), \(E|Y|\), \(E|XY|\) are finite, then
      \(E(XY\mid\mathcal H)=E(X\mid\mathcal H)\,E(Y\mid\mathcal H)\) almost
      surely.
\item \(\mathcal A\perp\mathcal B\mid\mathcal H\) holds exactly when
      \(P(A\mid\mathcal B\vee\mathcal H)=P(A\mid\mathcal H)\) almost surely
      for every \(A\in\mathcal A\).
\item \(\mathcal A\perp\mathcal B\mid\mathcal H\) holds exactly when, for
      every \(A\in\mathcal A\), some \(\mathcal H\)-measurable random
      variable is a version of \(P(A\mid\mathcal B\vee\mathcal H)\).
\item If \(\mathcal A\perp\mathcal B\mid\mathcal H\) and \(X\) is an
      integrable \(\mathcal A\)-measurable real random variable, then
      \(E(X\mid\mathcal B\vee\mathcal H)=E(X\mid\mathcal H)\) almost surely.
\item If \(\mathcal B\subseteq\mathcal H\), then
      \(\mathcal A\perp\mathcal B\mid\mathcal H\). If every event of
      \(\mathcal H\) has probability \(0\) or \(1\), then
      \(\mathcal A\perp\mathcal B\mid\mathcal H\) holds exactly when
      \(P(A\cap B)=P(A)P(B)\) for all \(A\in\mathcal A\), \(B\in\mathcal B\).
      If \(\mathcal H=\mathcal F_0\), then
      \(\mathcal A\perp\mathcal B\mid\mathcal H\) for all \(\mathcal A,\mathcal B\).
\item For random elements \(X,Y,Z\) with values in measurable spaces,
      \(X\perp Y\mid Z\) holds exactly when
      \(E[u(X)v(Y)\mid Z]=E[u(X)\mid Z]\,E[v(Y)\mid Z]\) almost surely for
      all bounded measurable real \(u,v\). For standard Borel spaces this is
      the relation of \cref{def:p16-measurable-ci}.
\end{enumerate}
\end{hbtheorem}

Items~1--5 are Lauritzen's Lemma~2.6, Theorem~2.9, Theorem~2.10,
Corollary~2.11 and Corollary~2.12, and item~6 collects his
Examples~2.2, 2.7 and~2.13
\cite[pp.~40--46]{Lauritzen2019GraphicalLectures2020}. For item~2 the
source proves indicators and then extends; the annex writes the extension
through simple, nonnegative and signed variables. Lauritzen reads
item~3 as saying that, once the information in \(\mathcal H\) is
available, \(\mathcal B\) adds nothing for predicting events of
\(\mathcal A\) (p.~45). Item~3 is one-sided in form although the relation
is symmetric; the inference rules below are proved from it.

\begin{hbtheorem}{thm:p15-w02c-p15-009-sigma-ci-rules}{proved=annex:p15-w02c-p15-009-sigma-ci-rules}{Symmetry, reduction, decomposition, weak union and contraction}
Let \(\mathcal A,\mathcal B,\mathcal C,\mathcal D\) be sub-sigma-algebras of
\(\mathcal F_0\).
\begin{enumerate}
\item \emph{Symmetry.} \(\mathcal A\perp\mathcal B\mid\mathcal C\) implies
      \(\mathcal B\perp\mathcal A\mid\mathcal C\).
\item \emph{Reduction.} \(\mathcal A\perp\mathcal B\mid\mathcal C\) and
      \(\mathcal D\subseteq\mathcal B\) imply
      \(\mathcal A\perp\mathcal D\mid\mathcal C\). In particular
      (\emph{decomposition}) \(\mathcal A\perp(\mathcal B\vee\mathcal D)\mid\mathcal C\)
      implies both \(\mathcal A\perp\mathcal B\mid\mathcal C\) and
      \(\mathcal A\perp\mathcal D\mid\mathcal C\).
\item \emph{Weak union.} \(\mathcal A\perp(\mathcal B\vee\mathcal C)\mid\mathcal D\)
      implies \(\mathcal A\perp\mathcal B\mid(\mathcal C\vee\mathcal D)\).
\item \emph{Contraction.} \(\mathcal A\perp\mathcal C\mid\mathcal B\) and
      \(\mathcal A\perp\mathcal D\mid(\mathcal B\vee\mathcal C)\) imply
      \(\mathcal A\perp(\mathcal C\vee\mathcal D)\mid\mathcal B\).
\item \emph{Random elements.} For random elements \(X,Y,Z,W\) with values
      in measurable spaces, \(\sigma((Y,Z))=\sigma(Y)\vee\sigma(Z)\), and for
      a measurable map \(\varphi\), \(\sigma(\varphi(Y))\subseteq\sigma(Y)\).
      Hence \(X\perp Y\mid Z\) implies \(Y\perp X\mid Z\) and
      \(X\perp\varphi(Y)\mid Z\); \(X\perp(Y,Z)\mid W\) implies
      \(X\perp Y\mid(Z,W)\); and \(X\perp Z\mid Y\) with
      \(X\perp W\mid(Y,Z)\) implies \(X\perp(Z,W)\mid Y\).
\end{enumerate}
\end{hbtheorem}

Items~1--4 are properties (A1)--(A4) of Lauritzen's Theorem~2.18 and
item~5 is his Corollary~2.20, properties (C1)--(C4)
\cite[pp.~48--49]{Lauritzen2019GraphicalLectures2020}. On a finite set of
labelled variables these rules make the relation
\(A\perp B\mid C\Leftrightarrow X_A\perp X_B\mid X_C\) a semi-graphoid in
the sense of Lauritzen's (S1)--(S4) (p.~52); the source reports that
Pearl (1988) conjectured (C1)--(C4) to be a complete axiomatization of
probabilistic conditional independence, that completeness fails, and that
Studen\'y (1992) showed that no finite axiomatization exists (p.~52). Neither
of those two works is cited here directly.

Theorem~2.18 of the source also lists an intersection property (A5), whose
conclusion conditions on \(\mathcal B\cap\mathcal C\) and is proved after
the null sets of \(P\) have been added to the sigma-algebras; Lauritzen
notes that without that augmentation the property is not true in general
(p.~48). Item~3 of the next theorem is his Example~2.19 (p.~49), and shows
the failure for the plain sigma-algebras of this chapter. For random
variables, the corresponding failure of the intersection rule (C5)
without positivity is item~2 of \cref{thm:p15-w02c-p15-009-undirected-gap}.

\begin{hbtheorem}{thm:p15-w02c-p15-009-ci-tables}{proved=annex:p15-w02c-p15-009-ci-tables}{Conditioning can create or destroy independence, and intersection fails}
Let \(C\) be an event with \(0<P(C)<1\) and
\(\mathcal H=\{\varnothing,C,C^c,\Omega\}\).
\begin{enumerate}
\item Suppose \(\Omega\) is partitioned by \(A,A^c\), \(B,B^c\), \(C,C^c\)
      into eight cells with probabilities, in eighteenths,
      \[
      \begin{array}{c|cc|cc}
        & \multicolumn{2}{c|}{\text{on }C} & \multicolumn{2}{c}{\text{on }C^c}\\
        & B & B^c & B & B^c\\\hline
       A & 2 & 1 & 2 & 4\\
       A^c & 4 & 2 & 1 & 2
      \end{array}
      \]
      Then \(A\perp B\mid\mathcal H\), but \(P(A\cap B)=2/9\ne1/4=P(A)P(B)\).
\item With probabilities, in twelfths,
      \[
      \begin{array}{c|cc|cc}
        & \multicolumn{2}{c|}{\text{on }C} & \multicolumn{2}{c}{\text{on }C^c}\\
        & B & B^c & B & B^c\\\hline
       A & 1 & 2 & 2 & 1\\
       A^c & 2 & 1 & 1 & 2
      \end{array}
      \]
      \(A\) and \(B\) are independent, but \(A\perp B\mid\mathcal H\) fails.
\item Let \(\Omega=\{0,1\}^3\), with all subsets measurable, and give the
      points \((0,0,0)\) and \((1,1,1)\) probability \(1/2\) each. Let \(X_i\) be the \(i\)-th
      coordinate and \(\mathcal A_i=\sigma(X_i)\). Then
      \(\mathcal A_1\perp\mathcal A_2\mid\mathcal A_3\) and
      \(\mathcal A_1\perp\mathcal A_3\mid\mathcal A_2\), and
      \(\mathcal A_2\cap\mathcal A_3=\{\varnothing,\Omega\}\), yet the
      relation
      \[
       \mathcal A_1\perp(\mathcal A_2\vee\mathcal A_3)\mid\mathcal A_2\cap\mathcal A_3
      \]
      fails.
\end{enumerate}
\end{hbtheorem}

The two tables are Lauritzen's Example~2.3, and he asserts both
conclusions with ``as can readily be checked''
\cite[p.~41]{Lauritzen2019GraphicalLectures2020}; the annex checks them.
Item~1 is the pattern in which a common cause, coded by \(C\), induces an
association that disappears within each value of \(C\); item~2 shows that
conditioning on a further event can make independent events dependent.
In item~3 the three coordinates agree almost surely, so each of
\(X_2,X_3\) carries all the information of the other and of \(X_1\), and
the two premises of an intersection rule hold while its conclusion fails.
Lauritzen adds that once the null sets are adjoined, the intersection of
the two sigma-algebras becomes \(\mathcal F_0\), so that the augmented
conclusion holds trivially (p.~49).
